%! Characteristics of Quadratic Functions — Unit 2, Lesson 2.0 — Student Packet Cover Sheet
% 12pt like every other student component. The banner is \coverbanner, which measures the title
% block and sizes the band to it. \namedateperiod appears HERE AND NOWHERE ELSE in the packet.
% Three scored rows: Warm-Up, Guided Notes & Practice, Homework — four cells in every row of the
% tocbox table (c l X r) or the widths collapse. Homework is due the first class after two study
% halls — course policy of record — never "next class". The Keep in Mind box is a CONTENT summary.
% ONE PAGE: four targets and a short Keep in Mind, matched to unit01/lesson04's weight (the first
% build came out 2pp on five long targets and a paragraph-length remind box).
\documentclass[12pt]{article}
\usepackage{algebra2-article}
\usepackage{algebra2-boxes}

\begin{document}

\coverbanner{Unit 2 \quad Quadratic Functions}{Lesson 2.0 \quad Characteristics of Quadratic Functions}

\namedateperiod
\vspace{0.18in}

\boxguard[14]
\begin{learningtargetbox}
\small
\begin{enumerate}[leftmargin=1.5em, itemsep=5pt]
  \item \ldots{} find the \textbf{vertex} of a \textbf{parabola} from its graph and say whether it
        is a \textbf{maximum} or a \textbf{minimum} point, using the direction it opens.
  \item \ldots{} state that \textbf{maximum} or \textbf{minimum value} --- an \emph{output}, with
        units --- and separately say \emph{where} it occurs, an \emph{input}.
  \item \ldots{} write the \textbf{axis of symmetry} $x=h$ and use symmetry to pair equal outputs
        and to locate the vertex from a table or from the two \textbf{zeros}.
  \item \ldots{} read the \textbf{zeros}, the \textbf{$y$-intercept}, the \textbf{increasing} and
        \textbf{decreasing} intervals, the \textbf{domain}, the \textbf{range}, and the
        \textbf{end behavior}.
\end{enumerate}
\end{learningtargetbox}

\vspace{0.14in}

\boxguard[14]
\begin{tocbox}
\small
\renewcommand{\arraystretch}{1.5}
\begin{tabularx}{\linewidth}{@{} c l X r @{}}
\rowcolor{forestbg}
\textbf{\#} & \textbf{Component} & \textbf{Description} & \textbf{Score} \\
\midrule
1 & Warm-Up & Evaluating a rule with negatives and squares, a linear graph's features, and one
              parabola with its high point circled & \blank{1.2cm} \\
2 & Guided Notes \& Practice & A ball off a 48-foot ledge: where the curve turns, how high that
              is, where it crosses, where it rises and falls, and what both ends do. Four rows of
              notes, then one graph read together & \blank{1.2cm} \\
3 & Homework & A graph read whole, a contrast pair, a table and its second differences, the
              special cases, a punt, and one multiple-choice item ---
              \textbf{scored; due the first class after two study halls} & \blank{1.2cm} \\
\midrule
  & \multicolumn{2}{r}{\textbf{Total}} & \blank{1.2cm} \\
\end{tabularx}
\end{tocbox}

\vspace{0.14in}

\boxguard[8]
\begin{remindbox}
\small
A \textbf{parabola} is the graph of $f(x)=ax^2+bx+c$, $a\ne0$: it opens \emph{up} when $a>0$ and
\emph{down} when $a<0$, and its \textbf{vertex} is the one point where it turns. \textbf{The
maximum or minimum \emph{value} is an output --- a $y$-value; \emph{where} it happens is an input,
the $x$-coordinate of the vertex.} The \textbf{axis of symmetry} $x=h$ is the mirror line through
the vertex, so inputs equally far either side have equal outputs --- which is why the two
\textbf{zeros} are symmetric about it and averaging them finds it. A parabola has \emph{two, one,
or no} zeros, never three; its \textbf{domain} is all real numbers, its \textbf{range} is $y\ge k$
or $y\le k$, and both ends of its graph always go the \emph{same} way.
\end{remindbox}

\end{document}
